\documentclass[12pt]{article}
\usepackage[utf8]{inputenc}
\usepackage[T1]{fontenc}
\usepackage{hyperref}
\title{Tirmaga-Chai Suri Language Resources for K-12 Curriculum Development}
\author{Compiled by Claude Sonnet 5.5 (claude-sonnet-5-5) for Andr\'as Kornai}
\date{}
\begin{document}
\maketitle

\section{Introduction}
Tirmaga-Chai Suri (ISO 639-3 \texttt{suq}) is a Surmic language of the West Omo Zone, Ethiopia. Its dialects are Tirma and Chai. Speaker counts in the sources vary: 27,000 (2007 census, per Wikipedia), 55,000 (Joshua Project, citing Ethnologue 2016) and 19,622 (Global Recordings Network). Search found few resources, and none of them are school textbooks.

\section{Primary Educational Resources}
A photographer's profile says the Suri Literacy Project ran for about 20 years from 2001 and produced drawings and stories for students. No primer or textbook was found in search results.
\begin{itemize}
\item LensCulture, ``The Eloquent Suri of Ethiopia'': \url{https://www.lensculture.com/projects/1642756-the-eloquent-suri-of-ethiopia}
\end{itemize}

\section{Linguistic and Reference Materials}
\begin{itemize}
\item Wikipedia, ``Suri language'' (lists Latin script and a bibliography): \url{https://en.wikipedia.org/wiki/Suri_language}
\item Verbix/Endangered Languages page on Suri. It says Ethiopic script was used until 2009 and Latin since: \url{https://docs.verbix.com/EndangeredLanguages/Suri}
\item Endangered Languages Project bibliography for Suri: \url{https://www.endangeredlanguages.com/lang/5057/bibliography}
\end{itemize}

\section{Religious and Community Resources}
\begin{itemize}
\item Joshua Project, Suri, Tirmaga-Chai in Ethiopia. It reports Bible portions but no complete New Testament or Bible: \url{https://joshuaproject.net/people_groups/15130}
\item Scripture Earth listing for \texttt{suq}: \url{https://scriptureearth.org/00eng.php?iso=suq}
\end{itemize}

\section{Digital and Audio Resources}
\begin{itemize}
\item Global Recordings Network, Suri Tirmaga-Chai (audio Bible stories and other recordings): \url{https://globalrecordings.net/en/language/suq}
\end{itemize}

\section{Curriculum Development Context}
The reference description says Suri is taught as a subject in some schools, but no syllabus or textbook was found online. The script changed from Ethiopic to Latin around 2009, so older materials may use a different orthography. Any curriculum work should be checked with the Suri community and local education authorities.

\end{document}